\bookchapter{第十一章：攝影與圖像管理}

\chapterintro{實言之，衆人以手機拍照，所出之效遠不如手機本身之潛力。非君技不行，乃諸多小設置與小習君未察也。此章不述虛者，惟言君可立用之物。}

\booksection{一、君是否亦犯此等錯？}

先言最常見之數患，君視中了幾條：

\begin{itemize}
  \item \textbf{鏡頭乃髒}：此誠最普遍之患。掏手機之時指紋按於鏡頭上，拍出霧濛濛，還以爲手機壞矣。養習，拍照前以衣角拭鏡頭。
  \item \textbf{從不調曝光}：默認之式下，稍點屏乃對焦，然若點畢上下滑，可調亮度。衆人所拍之照非暗即亮，即不知此操也。
  \item \textbf{閃光燈永開}：手機LED閃光燈所拍之效誠醜，光又硬又冷，把人拍如鬼。除非真伸手不見五指，否則關之。暗光之下，以夜景模式遠勝閃光燈。
  \item \textbf{放大賴兩指拉伸}：於屏幕上雙指縮放乃數碼變焦，本質就是裁切畫面而後放大，畫質將明顯下降。若欲拍遠處之，直接用取景框上之變焦按鈕（0.5x/1x/2x/3x）。其中 0.5x（超廣角）與 1x（主攝）乃原生光學鏡頭，不損失畫質；2x/3x 則多爲"光學品質變焦"——通過高像素傳感器裁切實現，畫質接近光學變焦但並非獨立鏡頭。超過 3x 之變焦通常就是混合變焦或數碼變焦矣，畫質將有明顯下降。
\end{itemize}

\booksection{二、數立升出片率之小術}

此等無需何理論知識，用一次即可見效：

\begin{itemize}
  \item \textbf{三分法構圖+網格線}：於相機設置中開"網格線"（亦稱參考線）。拍照時將君欲突出之物置四線之交點上，畫面立"對"甚多。此乃攝影中最簡亦最效之一規，無甚玄學，即合人眼之視覺習慣。
  \item \textbf{覓"引導線"}：拍風景或街景之時，留意畫面中有無自然之線——馬路、欄杆、牆壁、一行樹——使此等線自畫面邊緣延至主體。此術善用，路人觀君所拍將以爲君學過攝影。
  \item \textbf{勿老立而拍}：衆人拍照永爲己眼之高，故所出之視角千篇一律。試蹲下自低角度拍之，或覓高處俯拍，同地所出之感全異。拍小孩與寵物尤需蹲下，平視方能拍出自然之感。
  \item \textbf{光線乃第一要素}：午時大太陽下拍人像，面上之影又硬又醜。最佳之光乃清晨與傍晚之"黃金時段"，陽光斜、色偏暖，拍何皆佳。陰天之光雖不出彩，然勝在均勻柔和，拍人像反較大太陽下易出片。
  \item \textbf{夜景模式使用要領}：開夜景模式之後，手機自動連拍多張而後合成。此時最要者乃——手絕不可抖。能靠牆即靠牆，能放桌上即放桌上，實在不可即屏息。夜景模式需君穩住二至四秒，抖則畫面糊矣。
\end{itemize}

\booksection{三、人像模式與專業模式}

\begin{itemize}
  \item \textbf{人像模式}：能拍出背景模糊、主體清晰之效（淺景深）。手機因鏡頭物理尺寸小，真實虛化之能有限，故賴AI算法以成——簡言之即AI將君之人摳出，繼而模糊背景。用時保持一至二米之距效最佳，離太遠或太近算法皆易出錯。又，拍人像之時，注意背景中勿有太複雜之線，否則AI扣圖將於髮際邊緣"翻車"。
  \item \textbf{專業模式}：若君欲全控畫面之效，可試專業模式。數個關鍵參數當知：ISO越低畫質越淨，然畫面將暗；快門越慢進光越多，然手抖將糊；白平衡管色之冷暖，調不佳照將偏藍或偏黃。常人無需常年用專業模式，然解此數參數至少可使君知相機在爲何。
\end{itemize}

\booksection{四、後期修圖：勿過度，恰到好處即足}

衆人對修圖有誤解，以爲修即不真。實則相機所拍之照本與肉眼所見不同，適當調整乃爲還君當時所見之畫面，非造假也。

\begin{itemize}
  \item \textbf{Snapseed}：谷歌出之免費修圖工具，功能專業然無廣告。吾日常修圖之流乃：先以"調整圖片"拉亮度與對比度 以"突出細節"銳化之 以"白平衡"糾偏色 終"裁"成所欲之構圖。Snapseed之"局部"甚佳，可僅調畫面中某一區域，無需全圖同動。
  \item \textbf{醒圖/美圖秀秀}：此二者宜修人像。然吾以爲有一則：美顏濾鏡拉至三成內即佳，逾此數即假，膚如塑料。自然較完美更好看。
  \item \textbf{Lightroom手機版}：若君對色有追求，此乃最專業者。有曲線調整、HSL分色調節此等高級之能。然於多數人而言稍殺雞用牛刀。
  \item \textbf{修圖工作流之議}：吾己之習乃先拍好原片，歸家後於電腦屏或iPad上慢慢修，手機屏太小，色亦不準。若必於手機上修，記將屏亮度調至中偏上，勿於省電之式下修圖。
\end{itemize}

\booksection{五、照片太多奈何？整理之策}

手機相冊中堆數千張照片，欲覓一張特定者翻半日——此體驗衆人皆有。數實用之整理法：

\begin{itemize}
  \item \textbf{善用搜索}：今iPhone與大部分安卓手機之相冊皆支持智能搜索。君可搜"貓""海邊""身份證""花"此等關鍵詞，系統將以AI識別畫面內容助君覓出。又可搜日期與地點，如"去年五月""北京"。此能知者不多，然誠甚佳。
  \item \textbf{定期清理垃圾}：截圖、重複之照、拍糊之廢片，此等乃佔空間之元兇。iOS相冊有"重複項目"自動合併之能；安卓可往文件管理中搜"screenshots"文件夾，批量刪不需之截圖。吾每月花五分鐘清理，即可騰出數G。
  \item \textbf{建相簿，勿全賴默認排序}：系統默認按時排序無錯，然於"孩子之成長照""旅行照片""重要證件"此等跨時段之內容，君當手建相簿歸類。蘋果之"共享相簿"尚可與家人同往其中加照，一人建，全家皆可觀。
  \item \textbf{備份乃底線}：照片乃最不可替之數據——文件失可重建，聊天記錄無亦即無，然孩子之照、老人之視頻失即真無矣。必開雲備份。iCloud照片最省心，安卓這邊各品牌皆有己之雲服務，或以百度網盤/阿里雲盤自動備份。原則乃：至少有一份雲備份，勿將雞蛋置於同一籃中。
\end{itemize}

\booksection{六、截圖與錄屏}

\begin{itemize}
  \item \textbf{截圖}：同按電源鍵+音量加鍵，乃所有手機通用之法。多安卓手機尚支持三指下滑截圖、指關節雙擊截圖。長截圖（滾動截圖）於安卓上甚成熟——截畢一張後點"滾動"按鈕，手機自動往下翻頁並拼爲一張長圖，宜截聊天記錄與網頁。
  \item \textbf{錄屏}：自控制中心添"屏幕錄製"，稍點始錄，再點止。錄前長按錄製按鈕可擇是否錄麥克風之聲——君欲作教程解說即開麥，僅錄操作畫面即關之。錄畢之視頻將自動存至相冊。
\end{itemize}

\begin{tcbpractice}{練之：以手機成一次完整之拍照+修圖+整理}
\begin{enumerate}
  \item 開相機設置中之"網格線"，覓一有自然光之處拍三張照片（注意拭鏡頭、調曝光）。
  \item 擇最佳之一張，以Snapseed作基礎調整（亮度、對比度、裁剪）。
  \item 爲此照建一專屬相簿，名爲"吾之作品"。
  \item 爲此照建一專屬相簿，名爲"吾之作品"。
\end{enumerate}

\end{tcbpractice}
